\chapter{Introduction}
\label{chap:introduction}

% See context/outline.md Ch 1 for paragraph-level plan.

\section{Background and Motivation}
\label{sec:background}

Norwegian transport companies move large volumes of goods every day. In 2025, these companies carried about 252 million tonnes on behalf of their customers, almost all of it within Norway \parencite{ssb2026godstransport}. Behind every load, in every transport company, a person called the traffic coordinator decides who drives what, in which vehicle, and when. These are large decisions, made dozens of times a day, and the choices the coordinator makes go directly to the company's results. Assigning a driver who is already close to the weekly hours cap turns the next load into overtime, which lifts cost and cuts margin. Letting a vehicle and its driver wait between two assignments means the company keeps paying for them while they earn nothing. Concentrating the heaviest workload on the same few drivers week after week wears them down and lowers employee satisfaction. Each of these is a specific choice, and together they roll up to the three outcomes a transport-company owner watches: how much overtime is run, how much idle time sits between assignments, and how evenly the workload is shared across the team. How that choice is made today, and what tools support or constrain it, was explored through interviews with traffic coordinators in Norwegian transport companies and through an ongoing dialogue with Admmit, the Norwegian software firm that offered this thesis as a bachelor task. Across those conversations, the same pattern surfaced. The information needed to track overtime, idle time, and load balance is in the company's systems, but it is spread across an order tool, a timesheet, the coordinator's spreadsheets, and the coordinator's memory, and it is hard to pull together into a single working overview. The coordinator can answer any specific question about a specific driver but cannot easily step back and see how the week as a whole is shaping up for the team. The thesis takes this as its starting point: making the data already in the company's systems easy to read at a glance, and using that data to propose a balanced plan the coordinator can review and adjust.

The data the previous paragraph described is already in the company's systems, and the technology to plan that data automatically already exists; the open question is what shape such a system should take to be acceptable to the coordinator who would have to run it. Admmit, the Norwegian software firm that offered this thesis as a bachelor task, has worked with transport-company owners for years; their main objective for any such system is to improve resource utilization for their transport-company customers, on the grounds that better-balanced plans cut overtime and idle hours and raise the company's earnings. Two requirements translate that objective into something the day-to-day user can accept. The first is a Human-in-the-Loop (HITL) operating model, so the coordinator can inspect, modify, accept, or reject every algorithm-generated assignment before it takes effect. The second is a multi-tenant architecture, in which a single deployed system serves several customer companies at once with their data and rules kept separate.

The interviews with traffic coordinators put both requirements in front of the people who would actually run the system, and the answers came back consistent with them. Coordinators confirmed the requirements made sense and added an observation of their own that pointed in the same direction: today they have no overview of utilisation across drivers and vehicles, and they want to keep control of every assignment made on their behalf. Several said directly that they would welcome an algorithm proposing the day's assignments. Others were more cautious, pointing to the unpredictability of transport work as a reason to keep human judgement at the centre of every decision. No coordinator opposed an algorithm proposing assignments; what each insisted on was the right to override. The split between the firm and the owners articulating the demand for automation and the coordinators who must operate the result day to day is the structural pattern \textcite{bainbridge1983ironies} describes in \emph{Ironies of Automation}, and it runs across the empirical foundation of this thesis.

The reason the visibility gap and the open planning step have not been closed already is that current planning practice in Norwegian transport companies leans heavily on what the coordinator carries in their head, and the software available to them does not cover the planning step at all. Coordinators carry much of what is needed to plan well in their heads: who holds the right licence for which load, who knows the difficult routes, who is close to the overtime ceiling, which vehicle suits which job. This kind of knowledge, learned on the job and rarely written down, is called tacit knowledge, and it works only as long as the same coordinator stays in the role. The software they work alongside does not fill in for that knowledge. The two commercial Norwegian Transport Management Systems (TMS) named by interviewees, Timpex and Opter, handle order management and invoicing well, but neither generates assignment plans on its own, and Timpex is described as slow under heavy use. The other interviewed companies plan from internal tools, spreadsheets, memory, or the phone. None of the systems in current use takes on the assignment step itself: deciding which driver and vehicle perform which order. That is the gap a planning system would have to fill.

The gap a planning system would have to fill is, in one sense, what the coordinator is already doing in their head. The mental rules they apply, "this driver is close to overtime, give the load to that one", "this vehicle is too small for that pallet count, swap it", are themselves a kind of algorithm, but one that is never written down, never compared against alternatives, and never improved between shifts. Modern scheduling techniques can take that informal algorithm, make it explicit, and run it on the same data with more constraints, more candidate plans, and a faster search than any person can manage by hand. The opportunity is therefore not to replace the coordinator's reasoning but to give it a more capable computational form.

Ressursplanlegger (hereafter RP) is the artefact this thesis builds and evaluates around exactly that opportunity. RP is an algorithm-assisted planning platform: an algorithm proposes a plan for a given day, week, or longer horizon, and a drag-and-drop timeline lets the coordinator inspect, modify, accept, or reject every algorithm-generated assignment before it takes effect. Three different optimisation engines sit behind the proposal, each suited to a different point on the speed-and-quality trade-off, and the same artefact can be tuned per customer company through configurable preference weights without code changes. The work behind RP took those underlying algorithms, which on their own produce only a list of assignments, and built around them the timeline, the override controls, and the conflict view a coordinator needs to actually use the result. RP is not a Transport Management System in the sense Timpex and Opter are: it does not handle order management or invoicing. It addresses the planning step those systems leave open.


\section{Anchor Concepts}
\label{sec:anchors}

In the dialogue with Admmit and in the interviews with traffic coordinators, three things came up often enough to be useful as a guiding frame for the rest of the thesis. They are not the only criteria the system has to meet, but they are the ones the project returns to when deciding what to build, what to evaluate, and where to draw the boundaries: Efficiency, Trust/control, and Adaptability.

\subsection{Efficiency}
\label{sec:anchor-efficiency}

The system has to help the company make better use of its drivers and vehicles. Three things should improve: drivers should run less overtime, they should spend less time waiting between assignments, and the work should be shared more evenly across the team. Before the system can improve any of these, it first has to make them easy to see for the coordinator and for the owner. The data is already in the company, but no one can read it at a glance today.

\subsection{Trust/control}
\label{sec:anchor-trust-control}

The coordinator has to stay in charge of every assignment the algorithm proposes. For each one, the coordinator should be able to look at it, change it, accept it, or reject it before it takes effect. Trust in the algorithm grows over time, but only if the coordinator can see what it does and change it when needed. A system that made the plans on its own and left the coordinator with no way to step in would not work in the kind of company this thesis is built for.

\subsection{Adaptability}
\label{sec:anchor-adaptability}

The same system has to work for very different companies without being rebuilt each time. Transport companies vary a lot: a small operator with around ten drivers does not run the same way as a firm with forty-five vehicles split across three departments. The system has to fit both, with the rules and priorities adjusted per company in the settings rather than in the code itself.


\section{Research Questions}
\label{sec:research-questions}

\begin{quote}
How can an algorithm-assisted resource planning platform improve resource utilization in Norwegian transport companies (reducing overtime, idle time between assignments, and uneven driver load) while remaining accountable to the traffic coordinator who operates it?
\end{quote}

The main question pairs the two things the system has to deliver: a real improvement in how the company uses its drivers and vehicles, and a way of working that keeps the coordinator in charge of the result. A system that improved the numbers but took the decision away from the coordinator would not answer the question, and a system that left the coordinator fully in charge but did not improve anything would not answer it either.

Three sub-questions break the main question into smaller pieces that can be answered later in the thesis.

\begin{quote}
\textbf{SQ1}: To what extent and under which conditions does an algorithm-assisted planning platform reduce overtime, idle time, and load imbalance below the level produced by current Norwegian planning practice?
\end{quote}

SQ1 measures the system's effect on resource utilization. It connects to \textbf{Efficiency} and is answered in §5.1.1, drawing on the multi-engine benchmark on synthetic datasets that approximate the fleet ranges of the seven interviewed companies, with the comparison baseline taken from the interview accounts of how planning is done today.

\begin{quote}
\textbf{SQ2}: Through which interface and process mechanisms can the traffic coordinator inspect, modify, accept, or reject every algorithm-generated assignment without making the planning workflow heavier than current practice?
\end{quote}

SQ2 examines the specific mechanisms by which the coordinator stays in charge of every assignment the algorithm proposes. It connects to \textbf{Trust/control} and is answered in §5.1.2 by mapping the four required actions (inspect, modify, accept, reject) to the timeline, the override controls, and the conflict view in RP, against the human-in-the-loop theory established in \Cref{sec:hitl}.

\begin{quote}
\textbf{SQ3}: Through which configurable elements can the same deployed system function meaningfully across Norwegian transport companies that differ in operational rules, fleet size, and planning horizon, without changes to the code?
\end{quote}

SQ3 examines which parts of the system can be tuned per company at deployment and runtime. It connects to \textbf{Adaptability} and is answered in §5.1.3, drawing on the per-company configuration of the soft-constraint weights, the dynamic planning horizon set in the dashboard, and the multi-tenant architecture that keeps each company's data and rules isolated.


\section{Scope and Delimitations}
\label{sec:scope}

The system covers the planning step itself: turning a list of orders for a given period into a concrete plan that says who drives what, in which vehicle, and when. The length of that period is set by the coordinator in the dashboard and ranges from a single day to a week or longer, since the companies the system has to fit do not all plan on the same horizon. Around that core, the system holds the information the planning step needs: who the drivers are and what they are qualified to do, what vehicles the company has and what each one can carry, when people are working and when they are not.

The plan itself is produced by an automated solver. The system runs three different solver engines over the same input data, so a coordinator who needs a quick suggestion under time pressure and a coordinator who can wait for a more carefully searched plan are both served from the same artefact. Each engine fits a different point on the speed-and-quality trade-off, and \Cref{sec:iter-multiengine} describes what each one is and when each is preferred.

Different transport companies care about different things, and the priorities the solver should weigh need to be adjustable per company without changing code. The system exposes per-company configuration of the weights on the soft constraints (overtime cost relative to idle time, for example, or how strongly to balance workload across drivers), so a company that values one priority more highly than another can have that preference reflected in the plans the solver returns. Each of the three parts (the planning step, the solver, and the per-company weighting) is in scope because it serves at least one anchor: making utilisation visible and improving it under Efficiency, supporting override under Trust/control, and supporting per-company tuning under Adaptability.

Several capabilities are deliberately out of scope, and they are worth naming because they came up in the interviews and could otherwise be assumed by default. The system is not a driver-facing app and does not push notifications to drivers: it is a planning tool, not a dispatch tool. It does not handle billing or invoicing, which is what Timpex and Opter already cover, and bridging the two is a project of its own. It covers Norwegian transport only, since the regulatory frame is national. It does not track vehicles in real time on the road; it plans the day, it does not monitor the day's execution. It does not try to plan the order in which a single driver should visit a list of stops, sometimes called the Vehicle Routing Problem, since the question it answers is who does what, not in which order. And it does not let the algorithm run without a human in the loop; that is a deliberate Admmit requirement, not a setting that can be turned off.